\begin{abstract}
关于争议性议题的在线讨论日益受到协调性叙事攻击的干扰，而大语言模型使得此类攻击的规模化成本极低。现有研究主要聚焦于事后检测，对攻击如何在讨论中动态扎根的机制探索不足。为填补这一空白，我们提出\textbf{NarrateSim}——一个多智能体模拟框架，其中LLM智能体通过信念--意愿--意图（BDI）过程进行推理，并在无标度网络上交换消息，同时由提示驱动的攻击者注入策略性说服内容。实验揭示了两个相互关联的模式。在攻击端，即使少量攻击者也能使群体观点偏离无攻击基线；基于证据和权威凭证的策略最显著地重塑立场，而受害者式叙事则提升情绪温度；相同内容在不同受众中可能产生截然不同的效果。在防御端，\textbf{预防性接种}优于事后纠正，表明早期免疫比事后修复更为可靠。然而时机至关重要：一旦情绪升温，即使自我反思提示也可能放大而非削弱攻击效果。总体而言，本研究将叙事攻击研究从事后检测推进到动态机制分析，并为保护在线社区免受协调性操纵提供了可行的治理策略。
\end{abstract}
